% ==================================================================
% §15.5 OPENING — Introductory part of the cosmological ε-section
% Phase 2 Step 1 — FINAL (GPT round 14 applied: 5 corrections).
%
% Inserted between \label{sec:cosmo_constraints_rebuilt} and
% \subsubsection*{15.5.1 Structural definition ...}.
%
% GPT round 14 corrections:
%   (1) "arising from three-stage closure" → "motivated by ... tested
%       here at the effective level" (softens first-principles overclaim).
%   (2) "five largely independent empirical constraints" → "five
%       observationally distinct empirical constraints with different
%       kernels and different dominant systematics" (statistical
%       independence not established).
%   (3) Paragraph 3 gets an explicit question-framing sentence: "not
%       whether ECT predicts a rigorously constant drift parameter at
%       first principles, but whether a common effective ε-range can
%       account for several observationally distinct channels".
%   (4) Roadmap sharpens division of labour: synthesis in main text,
%       derivations in appendices.
%   (5) Closing sentence in §4 prevents "cherry-picking" reading:
%       excluded channels are explicitly discussed with reasons.
%
% NO legacy numbers. NO numerics (belong to §15.5.3). NO citations.
% ==================================================================

The previous part of Chapter~15 examined the separate cosmological
sectors --- flatness, horizon and monopole problems, age and lookback
integrals, structure-formation phenomenology --- within the
ordered-branch $\phi$-closure of ECT. The present section addresses a
different kind of question. In the standard phenomenological treatment
of late-time cosmological anomalies, each tension typically receives
its own parameterisation: a pre-recombination component for the Hubble
discrepancy, a modified gravity law for the late-time growth sector,
an abundance enhancement factor for the high-$z$ galaxy counts. ECT
offers a structurally different path. A single effective deformation
parameter, $\varepsilon$, motivated by the three-stage
condensate-closure picture and tested here at the effective level,
parametrizes the common late-time gravitational response function
$G_{\rm eff}(z)$ used across the cosmological channels considered
below. The section
therefore asks whether a single value of $\varepsilon$ --- or more
precisely, a single narrow range --- is compatible with a set of
observationally distinct channels probing that response.

If five channels of qualitatively different physics (a CMB-era
distance-ratio proxy, a high-$z$ structure-formation anomaly, direct
late-time Hubble-rate measurements, late-time growth likelihoods, and
late-time potential-evolution signals), when analysed separately
within the same effective framework, admit the same narrow
$\varepsilon$-interval, then that interval is not a fit of a
parameter to one phenomenon: it is the intersection of five
observationally distinct empirical constraints with different kernels
and different dominant systematics bearing on one common quantity. A
non-trivial intersection of this kind is a genuine indirect consistency
check of the framework. If the five channels had disagreed --- for
example, if any one of them had required a negative $\varepsilon$ that
the others excluded --- then the single-parameter description would
have failed in a clean, falsifiable way. Conversely, a narrow joint
band with five channels of different kernels and systematics
constitutes the strongest indirect consistency check available at
this methodological level that a single underlying mechanism may be
at work.

Throughout \S15.5 $\varepsilon$ is treated as an effective diagnostic
parameter rather than as a first-principles constant of ECT. The
retained-five-probe joint band reported in
\S\ref{subsec:eps_joint_band_rebuilt} is an empirical result under the
stated methodology, not a measurement of a fundamental theoretical
quantity and not a global statistical combination of independent
likelihoods. A first-principles closure-level derivation of the full
epoch-dependent function $\varepsilon(z)$ predicted by three-stage
condensate closure, together with its mapping to the effective
uniform-$\varepsilon$ layer used here, remains open. The question
pursued in \S15.5 is therefore not whether ECT predicts a rigorously
constant cosmological drift parameter at first principles, but whether
a common effective $\varepsilon$-range can account for several
observationally distinct channels within one diagnostic layer. In this
sense the retained band is compatible with $\varepsilon$ being a true
constant; it is equally compatible with a slow-varying $\varepsilon(z)$
averaged over the kernel of each probe. Distinguishing these two
possibilities belongs to the closure-level extension.

The section is organised as follows.
\S\ref{subsec:eps_def_rebuilt} fixes the structural meaning of
$\varepsilon$ within ECT and states the effective uniform-$\varepsilon$
ansatz. \S\ref{subsec:eps_retained_rebuilt} specifies the retained
five probes, what each of them actually constrains, and introduces the
two-sided vs one-sided-upper-bound distinction.
\S\ref{subsec:eps_joint_band_rebuilt} reports the joint effective
allowed band and fixes the working benchmark.
\S\ref{subsec:eps_excluded_rebuilt} discusses the channels not
retained, in three structurally distinct categories, together with the
reasons for their exclusion.
\S\ref{subsec:eps_comparison_rebuilt} places the result in the context
of alternative proposals for the same observational pattern.
\S\ref{subsec:eps_outlook_rebuilt} states what the analysis
establishes and what it does not, logs the first-principles open
problem \textbf{OP-Hubble-derive}, and outlines directions of future
work. The six subsections above give the synthesis and the cross-probe
reading; the detailed per-probe derivations, extraction logic,
uncertainty budgets, and the construction of each one-dimensional
$\varepsilon$-interval are to be provided in the dedicated methodology
appendices referenced from the subsections below as the rebuild is
completed. Channels not
retained in the joint band are still discussed explicitly in
\S\ref{subsec:eps_excluded_rebuilt}, together with the reasons why
they are methodologically limited, structurally outside the present
$\varepsilon$-sector, or not independent enough to serve as primary
constraints.

